\section{Preliminaries}
\label{sec:preliminaries}

We use $[a,b]=\{a,a+1,\ldots,b\}$ for integer intervals.  Let $\Sigma$ be a
finite alphabet of size $q$.  A nearest-neighbour $\ZZ^d$ shift of finite
type is specified by allowed symbol pairs on edges parallel to each standard
basis vector.  A pattern on a finite set is \emph{locally admissible} if all
nearest-neighbour edges contained in its support obey these rules.  For the
cube $Q_n=[1,n]^d$, write
\[
  C_n=\#\{\text{locally admissible patterns on }Q_n\}.
\]
For the presented nearest-neighbour system considered below, the topological
entropy is
\[
  h(X)=\lim_{n\to\infty}\frac{\log C_n}{n^d};
\]
in particular, $h(X)\le n^{-d}\log C_n$ for every $n$.  These standard
finite-volume formulations and their relation to entropy are discussed, for
example, in \cite{Friedland1997,LindMarcus2021}.

For $1\le i\le d$, let $R_i$ be the coordinate reflection on configurations,
\[
  (R_i x)_{(v_1,\ldots,v_i,\ldots,v_d)}
  =x_{(v_1,\ldots,-v_i,\ldots,v_d)}.
\]
The reflection hypothesis in this note is presentation-level: the chosen
nearest-neighbour rules are \emph{coordinate-reflection invariant} if every
allowed edge remains allowed when its orientation in the corresponding
coordinate is reversed.  We call a presented SFT coordinate-reflection
invariant when it is equipped with such a rule set.  This hypothesis implies
$R_iX=X$ for every $i$; the converse for an arbitrary redundant presentation
is neither asserted nor used.
For $m\ge1$, define
\[
 P_m=\bigl|\Fix(m\ZZ^d,X)\bigr|,
\]
where $\Fix(m\ZZ^d,X)$ consists of configurations invariant under translation
by $me_i$ for all $i$.  Restriction to $[1,m]^d$ is injective on this set and
its image is locally admissible, whence
\begin{equation}
 P_m\le C_m.
 \label{eq:periodic-upper-basic}
\end{equation}

For the one-dimensional statement, let $A$ be an $r\times r$ nonnegative
integral matrix and let $X_A$ be its edge shift.  Multiple edges are allowed.
The fixed-point count is
\begin{equation}
 F_n=|\Fix(\sigma^n,X_A)|=\tr(A^n).
 \label{eq:trace-fixed}
\end{equation}
Let $L_n$ be the number of points whose least positive period is exactly
$n$, and let $\tau(n)$ denote the number of positive divisors of $n$.
Decomposition by least period gives
\begin{equation}
 F_n=\sum_{e\mid n}L_e,
 \qquad
 L_n=\sum_{e\mid n}\mu(n/e)F_e,
 \label{eq:mobius}
\end{equation}
where $\mu$ is the M\"obius function.  Since every orbit of least period $n$
contains exactly $n$ points, the corresponding orbit count is $L_n/n$.
